%%%%%%%%%%  scan idx 224  |  folio 211  %%%%%%%%%%
(loc.cit.)
\[
  C^\infty(X,E)=\varprojlim_r C^r_2(X,E).
\]

\textbf{3.0.5.}
Soit $S$ un espace topologique. Nous aurons besoin de la notion de «~fibré hilbertien~» au-dessus de $S$, pour laquelle nous renvoyons le lecteur à [12] ou [14]. On a un foncteur additif de la catégorie des fibrés hilbertiens (réels resp. complexes) au-dessus de $S$ dans celle des $\cO_S$ rep. $\cO_S\otimes\mathbb C$-Modules, associant à chaque fibré hilbertien $E$ le faisceau $\underline E$ des sections continues de $E$.

\subsubsection*{3.1. Résultats techniques auxiliaires}
\textbf{Lemme 3.1.1.} Soient $X,Y$ des espaces topologiques, $X$ étant localement compact. Notons $\underline C'(X,Y)$ le préfaisceau défini sur la base des ouverts relativement compacts de $X$ par $\underline C'(X,Y)(U)=C(\bar U,Y)$ (3.0.1). Le morphisme canonique de restriction
\[
  \underline C'(X,Y)\longrightarrow \underline C(X,Y)
\]
induit un isomorphisme sur les faisceaux associés.

\emph{Preuve.} C'est immédiat.

\textbf{Lemme 3.1.2.} Soient $E,F$ des espaces vectoriels topologiques localement convexes (sur $\mathbb C$, pour fixer les idées), et $u:E\to F$ une application linéaire continue (resp. continue et d'image dense). Soit $X$ un espace compact. Munissons $C(X,E)$, $C(X,F)$ de la topologie de la convergence uniforme. Alors l'application linéaire
\[
  u_X:C(X,E)\longrightarrow C(X,F)
\]
induite par $u$ est continue (resp. continue et d'image dense).

